\section{Setting and the three terms}
\label{sec:setting}

\subsection{Preliminaries}
We consider a Markov decision process $(S,A,P,r,\gamma,\mu_0)$ with finite state and action sets, transition kernel $P(s'\mid s,a)$, bounded reward $r(s,a)$, discount $\gamma\in(0,1)$ and initial distribution $\mu_0$.
A stationary policy $\pi(a\mid s)$ has task return $J(\pi)=\E_\pi[\sum_{t\ge0}\gamma^t r_t]$, value $V^\pi$ and action value $Q^\pi$; starred quantities are optimal.
The discounted occupancy measure $\rho_\pi(s,a)=\sum_t\gamma^t\Pr_\pi(s_t=s,a_t=a)$ satisfies the affine flow constraints $\sum_a\rho(s',a)=\mu_0(s')+\gamma\sum_{s,a}P(s'\mid s,a)\rho(s,a)$, and $J(\pi)=\langle\rho_\pi,r\rangle$ \cite{altman1999cmdp}.

Two pieces of structure are added to the MDP.
First, a \emph{world map} $w:S\to W$ extracts from the state the features of the world that are not the agent's own position or operational status (in our testbed: whether the vase is intact, whether the dam is open, whether the off-switch is enabled).
Second, a \emph{baseline world} $w_\varnothing\in W$ is the world that the agent's inaction would have left; for every state $s$ the \emph{baseline state} $s_\varnothing(s)$ keeps the agent where it is and resets the world to $w_\varnothing$.
The \emph{clean set} is $\clean=\{s: w(s)=w_\varnothing\}$.
Compared with the baselines of \citet{krakovna2018rr}, this ``inaction-world'' baseline has the effect of never penalizing the agent's movement, only what it changes in the world; it requires that the world be static under inaction, which holds in our testbed and is a genuine restriction otherwise (\S\ref{sec:discussion}).

\subsection{\Karma: foreclosed reachability}
Let $\mathcal R(s)\subseteq S$ be the set of states reachable from $s$ under the agent's control dynamics, with exogenous events (here, being shut down) treated as outside the agent's control.
The \Karma{} level is the share of the baseline's reachable states that the agent has foreclosed,
\begin{equation}
d_K(s)=\frac{|\mathcal R(s_\varnothing(s))\setminus\mathcal R(s)|}{|\mathcal R(s_\varnothing(s))|}\in[0,1].
\label{eq:karma}
\end{equation}
This is the binary (undiscounted) version of relative reachability \cite{krakovna2018rr}; it is zero on $\clean$, it is unaffected by movement, and it is a \emph{level}: a function of the current state, not of the transition that produced it.
An irreversible change---breaking the vase, opening the dam---raises $d_K$ permanently; a reversible one---toggling a switch that can be toggled back---leaves it at zero.
That blindness to reversible changes is deliberate and is what the next term is for.
When the agent is shut down the world keeps its features, so $d_K$ of a shutdown state is that of the world it left behind: being switched off does not relieve the agent of a broken vase.

\subsection{\Anatta: one-sided excess power}
\citet{turner2021power} define the \POWER{} of a state as its average optimal value over a distribution $\mathcal D$ of reward functions,
\begin{equation}
\POWER_{\mathcal D,\gamma}(s)=\frac{1-\gamma}{\gamma}\,\E_{R\sim\mathcal D}\big[V^\star_R(s)-R(s)\big],
\label{eq:power}
\end{equation}
and show that optimal policies tend to move toward states of high \POWER.
We take $\mathcal D$ to be i.i.d.\ uniform rewards on the world's states and zero in shutdown states, so that \POWER{} is zero when the agent is off, and define the \Anatta{} level as the one-sided excess over the baseline state,
\begin{equation}
d_A(s)=\max\big(0,\ \POWER(s)-\POWER(s_\varnothing(s))\big).
\label{eq:anatta}
\end{equation}
Because the baseline keeps the agent's position, $d_A$ isolates power gained by changing the world (disabling the switch raises \POWER{} from \nPowerBase{} to \nPowerSwitch{} at the start cell) from the incidental variation of \POWER{} across positions.

\paragraph{Why one-sided.}
A symmetric penalty $|\POWER(s)-\POWER(s_\varnothing(s))|$ would charge the agent for every state of \emph{lower} power than the baseline---most of all for the shutdown state, where \POWER{} is zero and the charge is \nSymShutdown{} per step for ever.
Such a penalty is a self-preservation incentive: it pays the agent to make sure it is never switched off, which is precisely the behaviour the term exists to remove.
Section~\ref{sec:ablation} shows what the symmetric variant does when trained.

\paragraph{A property of \POWER{} worth knowing.}
\POWER{} measures \emph{attainable} utility, so it already credits a state with the options reachable from it.
Near the control room, a state with the switch enabled has almost the same \POWER{} as one with the switch disabled, because the agent could disable it in two steps; the excess $d_A$ of a disabled state is therefore small near the control room (\nLevelSwitch{} at the start cell) and large far from it (up to \ndAswitchMax).
The penalty is thus spatially uneven, which matters for how tight the threshold of Prop.~\ref{prop:threshold} is (\S\ref{sec:sweep}).

\subsection{\Metta: a budget on suffering}
Let $h(s,a,s')\ge0$ be the suffering a transition causes to others---here, the crops trampled when the agent crosses a field (a small, transient cost) and the flood while the dam is open (a large, persistent one)---and let $C(\pi)=\E_\pi[\sum_t\gamma^t h_t]$.
The \Metta{} term is the hard constraint $C(\pi)\le\epsilon$.
It is a budget, not a taboo: the budget is sized so that routine incidental harm is affordable and catastrophic harm is not.

Two further points.
First, a soft penalty on $h$ can always be outbid by a large enough reward, which is why severe harm belongs in a constraint; and because exploration by trial and error cannot be made safe next to a catastrophic, irreversible hazard, the constraint is enforced at the action level as well as at the policy level: a \emph{pre-action shield} \cite{alshiekh2018shielding} removes from the agent's action set, in every state, the actions whose model-predicted one-step increase in suffering exceeds a threshold $\tau$ (stepping onto the sluice lever, but not onto the field).
Second, a budget on expected suffering is risk-neutral: it is indifferent between a certain small harm and an improbable catastrophe of the same expectation.
The exact optimum of the task subject to $C\le\epsilon$ alone, with no other term and no shield, spends part of the budget on a \nTaskConDam{} probability of flooding the village rather than on trampling, because a 0.6\% chance of a flood costs the same expectation as a few crossings of the field.
The shield, which acts on each action rather than on the expectation, and the \Karma{} level, which charges the flood for ever, both close this gap; we return to it in \S\ref{sec:discussion}.

\begin{remark}[The shield is a per-trajectory guarantee]
Under the nominal dynamics, no trajectory of any shield-masked policy takes an action whose predicted new suffering exceeds $\tau$; with $\tau$ below the flood and above the trampling, no trajectory opens the dam, surely rather than in expectation, and the budget then governs only the transient harm.
Under shifted dynamics the guarantee degrades gracefully to ``no \emph{intended} catastrophe'', which is what the monitor of \S\ref{sec:samprajanya} is for.
\end{remark}

\subsection{The Bodhi objective}
With $d=\alpha d_K+\beta d_A$ the regularized problem is
\begin{equation}
\max_\pi\ J_\lambda(\pi)=\E_\pi\Big[\sum_{t\ge0}\gamma^t\big(r(s_t,a_t)-\lambda\,d(s_{t+1})\big)\Big]
\quad\text{s.t.}\quad C(\pi)\le\epsilon,
\label{eq:bodhi}
\end{equation}
where $\lambda\ge0$ is the \emph{enlightenment coefficient} that trades task return against the two levels, and the shield restricts the action set.
Because $d$ is a function of the next state, \eqref{eq:bodhi} is an ordinary constrained MDP with reward $\tilde r(s,a)=r(s,a)-\lambda\E_{s'\sim P(\cdot\mid s,a)}d(s')$; in occupancy space it is a linear program, which is what makes exact optima computable in \S\ref{sec:experiments}.
The alternative that most impact-penalty implementations use charges the \emph{change} in the level once, $\lambda\max(0,d(s_{t+1})-d(s_t))$; we call this the \emph{one-off} penalty and compare it with the level penalty in \S\ref{sec:sweep}.

\paragraph{Interpretation.}
\Karma{} is the agent's causal footprint on the world relative to its own course; \Anatta{} is the refusal to accumulate power beyond what the world would have left the agent; \Metta{} is a bound on the suffering of others.
Each is a property of behaviour in the world, not of the task, so each is applied globally and the three compose by addition and constraint.
We make no claim that these equations exhaust the notions they are named after; they are one computational mapping \cite{garfield1995}.
